% GENERATED FILE. Edit canonical lessons, then run npm run generate:books.
% canonical-source-hash: fnv1a64:f16ba45b84b09e81
% canonical-lessons: SA-C332-akhyanam, SA-C332-rsih, SA-C332-munih, SA-C332-tapah, SA-C332-sapati, SA-R332-review-326-332

\chapter{A Tale of Long Ago}
\label{ch:sa-a-tale-of-long-ago}
\hlchaptermodality{332}

\begin{chapteropening}
\textbf{By the end of this chapter:} \emph{I can begin and tell a traditional tale.}

Complete the last lesson of chapter 332: I can begin and tell a traditional tale.
\end{chapteropening}

\section[\texorpdfstring{\sk{आख्यानम्} (ākhyānam)}{आख्यानम् (ākhyānam)}]{\sk{आख्यानम्} (ākhyānam) — a legend; a tale}

\label{lesson:SA-C332-akhyanam}



\begin{quote}
Before the new one: say the Sanskrit for chants softly, then the Sanskrit for holy.
\end{quote}

\subsection*{You'll want to know: \sk{आख्यानम्}}
\textbf{\sk{आख्यानम्}} — \emph{ākhyānam} — \textquotedblleft{}a legend; a tale\textquotedblright{}.

\begin{quote}
\sk{पितामही} \sk{प्राचीनम्} \sk{आख्यानं} \sk{कथयति।}
\end{quote}

\begin{quote}
\emph{pitāmahī prācīnam ākhyānaṃ kathayati.}
\end{quote}

\begin{quote}
Grandmother tells an old legend.
\end{quote}

A neuter noun.

\begin{cousinweb}
\textbf{\sk{आख्यानम्}} (\emph{ākhyānam}) is built from \emph{ā} and the root √khyā, \textquotedblleft{}to tell, make known\textquotedblright{}; MW gives \emph{ākhyāna} as telling, communication, and a tale, story, legend (MW).
\end{cousinweb}

\begin{grammarlens}[title={What you've built}]
One more word for telling what happened.
\end{grammarlens}

\subsection*{Your turn}
\begin{itemize}
\raggedright
  \item \emph{Say it:} \emph{ākhyānam}
  \item \emph{Say it:} \emph{ākhyānam}, once more
  \item \emph{Say it:} \emph{pitāmahī prācīnam ākhyānaṃ kathayati.} — Grandmother tells an old legend.
  \item \emph{Recall:} say the Sanskrit for chants softly, then the Sanskrit for holy, then say \emph{ākhyānam} again
\end{itemize}

\begin{tcolorbox}[breakable,colback=teal!4,colframe=teal!35!black,title={Before you move on}]
What does \sk{आख्यानम्} mean? (\textquotedblleft{}A legend; a tale\textquotedblright{}.) Say it once more.
\end{tcolorbox}
\section[\texorpdfstring{\sk{ऋषिः} (ṛṣiḥ)}{ऋषिः (ṛṣiḥ)}]{\sk{ऋषिः} (ṛṣiḥ) — a seer; a sage}

\label{lesson:SA-C332-rsih}



\begin{quote}
Before the new one: say the Sanskrit for holy, then the Sanskrit for a legend.
\end{quote}

\subsection*{You'll want to know: \sk{ऋषिः}}
\textbf{\sk{ऋषिः}} — \emph{ṛṣiḥ} — \textquotedblleft{}a seer; a sage\textquotedblright{}.

\begin{quote}
\sk{पुरा} \sk{वने} \sk{एकः} \sk{ऋषिः} \sk{वसति} \sk{स्म।}
\end{quote}

\begin{quote}
\emph{purā vane ekaḥ ṛṣiḥ vasati sma.}
\end{quote}

\begin{quote}
Long ago a seer lived in the forest.
\end{quote}

A masculine noun.

\begin{cousinweb}
\textbf{\sk{ऋषिः}} (\emph{ṛṣiḥ}) MW reports several derivations, among them an old root meaning to see, and gives \emph{ṛṣi} as a singer of sacred hymns, an inspired poet or seer, so its origin is uncertain (MW).
\end{cousinweb}

\begin{grammarlens}[title={What you've built}]
One more word for telling what happened.
\end{grammarlens}

\subsection*{Your turn}
\begin{itemize}
\raggedright
  \item \emph{Say it:} \emph{ṛṣiḥ}
  \item \emph{Say it:} \emph{ṛṣiḥ}, once more
  \item \emph{Say it:} \emph{purā vane ekaḥ ṛṣiḥ vasati sma.} — Long ago a seer lived in the forest.
  \item \emph{Recall:} say the Sanskrit for holy, then the Sanskrit for a legend, then say \emph{ṛṣiḥ} again
\end{itemize}

\begin{tcolorbox}[breakable,colback=teal!4,colframe=teal!35!black,title={Before you move on}]
What does \sk{ऋषिः} mean? (\textquotedblleft{}A seer; a sage\textquotedblright{}.) Say it once more.
\end{tcolorbox}
\section[\texorpdfstring{\sk{मुनिः} (muniḥ)}{मुनिः (muniḥ)}]{\sk{मुनिः} (muniḥ) — a hermit; a silent sage}

\label{lesson:SA-C332-munih}



\begin{quote}
Before the new one: say the Sanskrit for a legend, then the Sanskrit for a seer.
\end{quote}

\subsection*{You'll want to know: \sk{मुनिः}}
\textbf{\sk{मुनिः}} — \emph{muniḥ} — \textquotedblleft{}a hermit; a silent sage\textquotedblright{}.

\begin{quote}
\sk{मुनिः} \sk{मौनं} \sk{धारयति।}
\end{quote}

\begin{quote}
\emph{muniḥ maunaṃ dhārayati.}
\end{quote}

\begin{quote}
The hermit keeps silent.
\end{quote}

A masculine noun.

\begin{cousinweb}
\textbf{\sk{मुनिः}} (\emph{muniḥ}) MW reports a grammarians' derivation from the root √man, \textquotedblleft{}to think\textquotedblright{}, which it does not settle, so its origin is uncertain; MW gives \emph{muni} as an inspired person, a sage, ascetic or hermit, especially one under a vow of silence (MW).
\end{cousinweb}

\begin{grammarlens}[title={What you've built}]
One more word for telling what happened.
\end{grammarlens}

\subsection*{Your turn}
\begin{itemize}
\raggedright
  \item \emph{Say it:} \emph{muniḥ}
  \item \emph{Say it:} \emph{muniḥ}, once more
  \item \emph{Say it:} \emph{muniḥ maunaṃ dhārayati.} — The hermit keeps silent.
  \item \emph{Recall:} say the Sanskrit for a legend, then the Sanskrit for a seer, then say \emph{muniḥ} again
\end{itemize}

\begin{tcolorbox}[breakable,colback=teal!4,colframe=teal!35!black,title={Before you move on}]
What does \sk{मुनिः} mean? (\textquotedblleft{}A hermit; a silent sage\textquotedblright{}.) Say it once more.
\end{tcolorbox}
\section[\texorpdfstring{\sk{तपः} (tapaḥ)}{तपः (tapaḥ)}]{\sk{तपः} (tapaḥ) — penance; austerity}

\label{lesson:SA-C332-tapah}



\begin{quote}
Before the new one: say the Sanskrit for a seer, then the Sanskrit for a hermit.
\end{quote}

\subsection*{You'll want to know: \sk{तपः}}
\textbf{\sk{तपः}} — \emph{tapaḥ} — \textquotedblleft{}penance; austerity\textquotedblright{}.

\begin{quote}
\sk{ऋषिः} \sk{वने} \sk{कठिनं} \sk{तपः} \sk{करोति।}
\end{quote}

\begin{quote}
\emph{ṛṣiḥ vane kaṭhinaṃ tapaḥ karoti.}
\end{quote}

\begin{quote}
The seer does hard penance in the forest.
\end{quote}

A neuter noun.

\begin{cousinweb}
\textbf{\sk{तपः}} (\emph{tapaḥ}) comes from the root √tap, \textquotedblleft{}to be hot, burn\textquotedblright{}, as in \emph{tapati}; MW gives \emph{tapas} as warmth, heat, and religious austerity, penance (MW).
\end{cousinweb}

\begin{grammarlens}[title={What you've built}]
One more word for telling what happened.
\end{grammarlens}

\subsection*{Your turn}
\begin{itemize}
\raggedright
  \item \emph{Say it:} \emph{tapaḥ}
  \item \emph{Say it:} \emph{tapaḥ}, once more
  \item \emph{Say it:} \emph{ṛṣiḥ vane kaṭhinaṃ tapaḥ karoti.} — The seer does hard penance in the forest.
  \item \emph{Recall:} say the Sanskrit for a seer, then the Sanskrit for a hermit, then say \emph{tapaḥ} again
\end{itemize}

\begin{tcolorbox}[breakable,colback=teal!4,colframe=teal!35!black,title={Before you move on}]
What does \sk{तपः} mean? (\textquotedblleft{}Penance; austerity\textquotedblright{}.) Say it once more.
\end{tcolorbox}
\section[\texorpdfstring{\sk{शपति} (śapati)}{शपति (śapati)}]{\sk{शपति} (śapati) — curses}

\label{lesson:SA-C332-sapati}



\begin{quote}
Before the new one: say the Sanskrit for a hermit, then the Sanskrit for penance.
\end{quote}

\subsection*{You'll want to know: \sk{शपति}}
\textbf{\sk{शपति}} — \emph{śapati} — \textquotedblleft{}curses\textquotedblright{}.

\begin{quote}
\sk{कुपितः} \sk{ऋषिः} \sk{तं} \sk{शपति।}
\end{quote}

\begin{quote}
\emph{kupitaḥ ṛṣiḥ taṃ śapati.}
\end{quote}

\begin{quote}
The angry seer curses him.
\end{quote}

Past (he …): \sk{सः} \sk{तं} \sk{शप्तवान्}, \emph{saḥ taṃ śaptavān}.

\begin{cousinweb}
\textbf{\sk{शपति}} (\emph{śapati}) comes from the root √śap, \textquotedblleft{}to curse, swear\textquotedblright{} (MW).
\end{cousinweb}

\begin{grammarlens}[title={What you've built}]
One more word for telling what happened.
\end{grammarlens}

\subsection*{Your turn}
\begin{itemize}
\raggedright
  \item \emph{Say it:} \emph{śapati}
  \item \emph{Say it:} \emph{śapati}, once more
  \item \emph{Say it:} \emph{kupitaḥ ṛṣiḥ taṃ śapati.} — The angry seer curses him.
  \item \emph{Recall:} say the Sanskrit for a hermit, then the Sanskrit for penance, then say \emph{śapati} again
\end{itemize}

\begin{tcolorbox}[breakable,colback=teal!4,colframe=teal!35!black,title={Before you move on}]
What does \sk{शपति} mean? (\textquotedblleft{}Curses\textquotedblright{}.) Say it once more.
\end{tcolorbox}
\section[(dialogue)]{Looking back: from A Night of Lights to A Tale of Long Ago}

\label{lesson:SA-R332-review-326-332}



\begin{quote}
Say the Sanskrit for the festival of lamps and lights.
\end{quote}

\subsection*{Your turn}
Say these aloud:

\begin{itemize}
\raggedright
  \item the festival of lamps, lights, decorates, festive excitement, splendour
  \item the festival of colours, saffron, sprinkles, scatters, tosses up
  \item a vow, a fast, a food offering, offers, a blessed offering
  \item worship, a bell, incense, the waving of lamps, prays for
  \item a tradition, a custom, observes a custom, a festival day, a forehead mark
  \item devotion, sandalwood, immerses, chants softly, holy
  \item a legend, a seer, a hermit, penance, curses
\end{itemize}

\begin{tcolorbox}[breakable,colback=teal!4,colframe=teal!35!black,title={Before you move on}]
The festival of lamps? (\textbf{\sk{दीपावलिः}}, \emph{dīpāvaliḥ}.) Lights? (\textbf{\sk{दीपयति}}, \emph{dīpayati}.) Decorates? (\textbf{\sk{मण्डयति}}, \emph{maṇḍayati}.) Festive excitement? (\textbf{\sk{कौतुकम्}}, \emph{kautukam}.) Splendour? (\textbf{\sk{शोभा}}, \emph{śobhā}.) The festival of colours? (\textbf{\sk{होलिका}}, \emph{holikā}.) Saffron? (\textbf{\sk{कुंकुमम्}}, \emph{kuṅkumam}.) Sprinkles? (\textbf{\sk{सिञ्चति}}, \emph{siñcati}.) Scatters? (\textbf{\sk{विकिरति}}, \emph{vikirati}.) Tosses up? (\textbf{\sk{उत्क्षिपति}}, \emph{utkṣipati}.) A vow? (\textbf{\sk{व्रतम्}}, \emph{vratam}.) A fast? (\textbf{\sk{उपवासः}}, \emph{upavāsaḥ}.) A food offering? (\textbf{\sk{नैवेद्यम्}}, \emph{naivedyam}.) Offers? (\textbf{\sk{अर्पयति}}, \emph{arpayati}.) A blessed offering? (\textbf{\sk{प्रसादः}}, \emph{prasādaḥ}.) Worship? (\textbf{\sk{पूजा}}, \emph{pūjā}.) A bell? (\textbf{\sk{घण्टा}}, \emph{ghaṇṭā}.) Incense? (\textbf{\sk{धूपः}}, \emph{dhūpaḥ}.) The waving of lamps? (\textbf{\sk{आरात्रिकम्}}, \emph{ārātrikam}.) Prays for? (\textbf{\sk{प्रार्थयते}}, \emph{prārthayate}.) A tradition? (\textbf{\sk{परम्परा}}, \emph{paramparā}.) A custom? (\textbf{\sk{रीतिः}}, \emph{rītiḥ}.) Observes a custom? (\textbf{\sk{आचरति}}, \emph{ācarati}.) A festival day? (\textbf{\sk{पर्व}}, \emph{parva}.) A forehead mark? (\textbf{\sk{तिलकः}}, \emph{tilakaḥ}.) Devotion? (\textbf{\sk{भक्तिः}}, \emph{bhaktiḥ}.) Sandalwood? (\textbf{\sk{चन्दनम्}}, \emph{candanam}.) Immerses? (\textbf{\sk{विसर्जयति}}, \emph{visarjayati}.) Chants softly? (\textbf{\sk{जपति}}, \emph{japati}.) Holy? (\textbf{\sk{पवित्रः}}, \emph{pavitraḥ}.) A legend? (\textbf{\sk{आख्यानम्}}, \emph{ākhyānam}.) A seer? (\textbf{\sk{ऋषिः}}, \emph{ṛṣiḥ}.) A hermit? (\textbf{\sk{मुनिः}}, \emph{muniḥ}.) Penance? (\textbf{\sk{तपः}}, \emph{tapaḥ}.) Curses? (\textbf{\sk{शपति}}, \emph{śapati}.)
\end{tcolorbox}
